\documentclass[exercice,a4]{thedocument}%
\usepackage{texmaths-tikz}
\usepackage{mdframed}
\startdatakeys
\source{25GENMATNC1}
\cycle{4}
\competencesDuSocle{CA2,RA1,CO2}
\connaissancesRequises{D2f*,D2e*}
\competencesTravaillees{D2e*,B3e*}
\enddatakeys
\begin{document}%
    Thomas souhaite construire le cerf-volant représenté par la figure ci-dessous~:\par\noindent
    \begingroup\centering
    \begin{tikzpicture}[baseline=(current bounding box.center)]
        \coordinate[label=above right:\mathspoint{E}](E)at(0,0);
        \coordinate[label=below:\mathspoint{B}](B)at(0,-3.46);
        \coordinate[label=above:\mathspoint{A}](A)at(0,1.54);
        \coordinate[label=right:\mathspoint{D}](D)at(2,0);
        \coordinate[label=left:\mathspoint{C}](C)at(-2,0);
        \fill[\currentcolor!50](A)--(C)--(B)--(D)--cycle;
        \tkzMarkRightAngle[size=0.5, colored](D,E,B)
        \tkzMarkSegments[mark=s||,colored](C,E E,D)
        \tkzFillAngle[\markscolor](D,B,E)
        \tkzLabelAngle[pos=1.25](D,B,E){30$^\circ$}
        \draw(A)--(C)--(B)--(D)--cycle;
        \draw(A)--(B);\draw(C)--(D);
    \end{tikzpicture}\hskip30pt
    \begin{minipage}{.2\linewidth}
        \begin{mdframed}
            On donne~:\par\vskip20pt\noindent
            $\widehat{\mathspoint{DEB}}=90^\circ$\par\noindent
            $\widehat{\mathspoint{EBD}}=30^\circ$\par\noindent
            $\mathspoint{AB}=50$ cm\par\noindent
            $\mathspoint{CD}=40$ cm\par\noindent
            $\mathspoint{ED}=\mathspoint{EC}$
        \end{mdframed}
    \end{minipage}\par\endgroup
    \begin{enumerate}
        \item Calculer \mathspoint{BE}. On donnera une valeur arrondie au
        millimètre.\par\noindent
        \begingroup\leavevmode\hspace*{-\leftmargin}%
        \begin{minipage}[c]{.6\linewidth}
            Lorsque Thomas a essayé son cerf-volant, il s'est demandé à quelle
            altitude il volait.\par\noindent Il a attaché sa corde à un piquet
            planté dans le sol (point \mathspoint{S}) puis est allé se placer
            (point \mathspoint{T}) parfaitement à la verticale sous son
            cerf-volant (point \mathspoint{H}).\par\noindent Il a alors mesuré
            certaines longueurs et a réalisé le schéma ci-contre.
        \end{minipage}\endgroup\hskip20pt
        \begin{tikzpicture}[scale=1.25, line width=1pt, baseline=(current bounding box.center)]
            \coordinate[label=below left:\mathspoint{S}](S)at(0,0);
            \coordinate[label=below right:\mathspoint{T}](T)at(7.60/4,0);
            \coordinate[label=above:\mathspoint{H}](H)at(7.60/4,19.04/4);
            \draw[fill=\currentcolor!50](S)--(T)--(H)--cycle;
            \path(S)++(158.23:0.4)coordinate(S1);
            \path(H)++(158.23:0.4)coordinate(H1);
            \draw[latex-latex](S1)--(H1);
            \path(S1)--(H1)node[midway,fill=white]{20,50 m};
            \path(S)++(0,-0.4)coordinate(S2);
            \path(T)++(0,-0.4)coordinate(T2);
            \draw[latex-latex](S2)--(T2);
            \path(S2)--(T2)node[midway,fill=white]{7,60 m};
        \end{tikzpicture}
        \item Calculer \mathspoint{HT}, altitude à laquelle volait son cerf-volant.
        \par\noindent On donnera une valeur arrondie au mètre.\par\bigskip
        \begingroup\leavevmode\hspace*{-\leftmargin}%
        \begin{minipage}{.6\linewidth}
            Il est conseillé de ne pas utiliser ce cerf-volant lorsque  le vent dépasse
            20 km/h.\par
            La météo annonce un vent ne dépassant pas 15 n\oe uds.\par\noindent
            On donne 1 n\oe ud = 0,514 m/s.
        \end{minipage}\endgroup
        \item Thomas peut-il faire voler son cerf-volant sans risque dans ces
        conditions~?
    \end{enumerate}
    \showthesource
\end{document}